% ID: SEL-D0021-OB-E02
% ESTADO: validada. S022, revisada con los criterios oficiales.
\textbf{a)} La derivabilidad exige continuidad. En $0$ el límite
izquierdo vale $-a$ y el valor de la función es $-1$, así que $a=1$.
Con continuidad asegurada, las pendientes laterales son el límite
de $-a/(x-1)^2$ por la izquierda y de $2x-b$ por la derecha.
Igualarlas da $-1=-b$, de modo que $\boxed{a=b=1}$.

\textbf{b)} Ahora $b=2$ y el punto está en el segundo tramo.
Como $f(2)=-1$ y $f'(2)=2$, la fórmula punto-pendiente da
$y+1=2(x-2)$, es decir, $\boxed{y=2x-5}$.
